% =====================================================================
% Mirante dos Dados — Working Paper #3
% RAIS, FAIRness e Lakehouse: replicação e extensão
% Autor: Leonardo Chalhoub  ·  Abril/2026
%
% NOTA: este é um SKELETON inicial. Será expandido com tabelas e figuras
% reais após a primeira execução do pipeline RAIS no Databricks.
% Spec completa em docs/vertical-rais-fair-lakehouse-spec.md
% =====================================================================
\documentclass[12pt,a4paper,oneside]{article}
\usepackage[T1]{fontenc}
\usepackage[utf8]{inputenc}
\usepackage[brazilian]{babel}
\usepackage{lmodern}
\usepackage[a4paper,top=3cm,bottom=2cm,left=3cm,right=2cm]{geometry}
\usepackage{setspace}
\usepackage{indentfirst}
\usepackage{titlesec}
\usepackage{caption}
\usepackage{booktabs}
\usepackage{graphicx}
\usepackage[hidelinks]{hyperref}
\usepackage{url}
\usepackage{float}

\onehalfspacing
\setlength{\parindent}{1.25cm}
\setlength{\parskip}{0pt}
\titleformat{\section}{\normalfont\bfseries\large\uppercase}{\thesection}{0.6em}{}
\titleformat{\subsection}{\normalfont\bfseries\normalsize}{\thesubsection}{0.6em}{}

\newcommand{\rs}{R\$\,}

\begin{document}

\begin{titlepage}
  \begin{center}
    \vspace*{1.5cm}
    {\small\textsc{Mirante dos Dados \quad\textbar\quad Working Paper n. 3}}\\
    {\small Versão 0.1 (skeleton) \textemdash{} Abril de 2026}

    \vfill

    {\LARGE\bfseries
     RAIS, FAIRNESS E LAKEHOUSE: REPLICAÇÃO E EXTENSÃO DE COMPARAÇÃO
     EMPÍRICA DE FORMATOS PARA BIG DATA PÚBLICO BRASILEIRO
     (2020\textendash 2025)}

    \vspace{1cm}
    {\large\itshape
     RAIS, FAIRness and Lakehouse: empirical replication and extension
     of format comparison for Brazilian public Big Data (2020\textendash 2025)}

    \vfill

    {\large\textbf{Leonardo Chalhoub}}\\[4pt]
    {\small Pesquisador independente \textemdash{} Engenharia de dados aplicada à transparência fiscal}\\[2pt]
    {\small\href{mailto:leonardochalhoub@gmail.com}{leonardochalhoub@gmail.com}}\\
    {\small\href{https://github.com/leonardochalhoub}{github.com/leonardochalhoub}}

    \vfill

    {\small Brasil \\ Abril de 2026}
  \end{center}
\end{titlepage}

\section*{Resumo}
Este artigo replica e estende a monografia de especialização do autor
\textit{"O Papel dos Metadados em Arquiteturas de Nuvem: FAIRness e
Governança de Dados em Big Data"} (UFRJ MBA Engenharia de Dados, 2023,
não publicada), atualizando-a com (i) instrumentos consagrados de
mensuração de FAIRness (FAIR Data Maturity Model do Research Data
Alliance), (ii) ampliação do conjunto de formatos comparados,
incluindo Apache Iceberg e Apache Hudi além de CSV, Apache Parquet e
Delta Lake, (iii) controle estatístico de variância nas medições de
desempenho (média, mediana, p95, IC 95\%) sobre múltiplas configurações
de cluster, e (iv) discussão crítica balanceada incluindo
\textit{when-not-to-use} para arquiteturas Lakehouse. O dataset é o
mesmo da monografia original — RAIS Vínculos Públicos
(\textasciitilde{}136 milhões de linhas, \textasciitilde{}62 GB
brutos), processado em arquitetura medallion sobre Databricks Free
Edition.

\vspace{0.4cm}\noindent\textbf{Palavras-chave:}
RAIS; FAIR principles; Lakehouse; Delta Lake; Apache Iceberg; Apache
Hudi; engenharia de dados; reprodutibilidade.

\section*{Status deste documento}
\textbf{Este é o esqueleto inicial} (versão 0.1). O conteúdo
substantivo (resultados empíricos, tabelas, figuras, discussão)
será expandido após a primeira execução do pipeline RAIS no
Databricks. Vide \texttt{docs/vertical-rais-fair-lakehouse-spec.md}
no repositório do projeto para a especificação detalhada das
extensões previstas.

\section{Introdução}\label{sec:intro}
[A ser escrito após a 1ª execução do pipeline.]

\section{Referencial Teórico}\label{sec:teoria}
[A ser escrito.]

\section{Metodologia}\label{sec:metodo}
[A ser escrito.]

\section{Resultados}\label{sec:resultados}
[A ser escrito após coleta empírica.]

\section{Discussão}\label{sec:discussao}
[A ser escrito.]

\section{Considerações Finais}\label{sec:conclusao}
[A ser escrito.]

\section*{Referências}
\addcontentsline{toc}{section}{Referências}
\begin{singlespace}\small
\noindent ARMBRUST, M. et al. Lakehouse: a new generation of open
platforms... \textbf{CIDR}, 2021.\par\smallskip
\noindent KUKREJA, M. \textit{Data Engineering with Apache Spark,
Delta Lake, and Lakehouse}. Packt, 2021.\par\smallskip
\noindent REIS, J.; HOUSLEY, M. \textit{Fundamentals of Data
Engineering}. O'Reilly, 2022.\par\smallskip
\noindent WILKINSON, M. D. et al. The FAIR Guiding Principles...
\textbf{Scientific Data}, v. 3, 2016.\par\smallskip
\noindent ZAHARIA, M. et al. Apache Spark: a unified engine for big
data processing. \textbf{CACM}, v. 59, n. 11, 2016.
\end{singlespace}

\end{document}
